
The comparative analysis and subsequent mechanistic investigation of the simulation results yielded a definitive conclusion regarding the optimization of solid waste management policies in the Municipality of Bacolod. The primary finding was that the conventional ``Status Quo'' approach, characterized by an equitable distribution of resources across all barangays with a heavy emphasis on Information, Education, and Communication (IEC), was mathematically insufficient to achieve municipality-wide compliance. While this strategy was capable of maintaining order in smaller, socially cohesive communities, it consistently failed to overcome the high behavioral friction present in densely populated urban centers. This failure was not due to a lack of total funding, but rather a strategic error in resource dilution; spreading the budget thinly prevented the Local Government Unit (LGU) from ever reaching the critical enforcement intensity required to break the inertia of non-compliance in high-resistance zones.

Furthermore, the simulation demonstrated that single-instrument policies---whether purely punitive or purely incentive-based---were equally flawed when applied uniformly. Pure enforcement collapsed due to its inability to sustain positive attitudes, while pure incentives failed due to the diminishing per-capita value of rewards in large populations. In contrast, the Deep Reinforcement Learning (DRL) agent's ``Sequential Saturation'' strategy provided a robust solution to this dilemma. By abandoning the heuristic of simultaneous equality in favor of sequential efficacy, the agent successfully leveraged the ``Graduation Effect,'' utilizing social norms to sustain compliance in stabilized areas while concentrating capital to besiege and conquer the most resistant urban targets. This confirmed that the optimal policy for a resource-constrained LGU was dynamic and adaptive: a phased intervention that treated compliance not as a static maintenance task, but as a series of strategic tipping points to be systematically triggered.

\section{Summary}
The evaluation of the Agent-Based Model through Global Sensitivity Analysis (GSA) and behavioral determinant mapping provided a robust computational foundation for the DRL agent's optimized policies. By transitioning from a ``Status Quo'' approach to a ``Sequential Saturation'' strategy, the model demonstrated how municipal resources could be leveraged more effectively by targeting the underlying mechanics of household decision-making \citep{ZhengAI2022}. 

These findings suggested that successful waste management policy in the local context was less about shifting public opinion and more about strategically reducing the structural friction inherent in the waste disposal process. The primary discovery of this evaluation was the dominance of convenience as a behavioral driver. With the Cost of Effort ($c_{\text{effort}}$) accounting for 83\% of output variance, the study confirmed that logistical friction was the single greatest inhibitor of waste segregation in the Municipality of Bacolod \citep{Yazawa2025}. Policies that did not actively reduce this ``cost'' were statistically likely to fail regardless of the level of public support.

Furthermore, the results highlighted a distinct ``Value-Action Gap'' characterized by the near-zero sensitivity of the Attitude ($w_a$) parameter. This mathematically illustrated that while IEC campaigns were foundational, they had reached a point of diminishing returns \citep{Zhao2022}. The DRL agent's strategic shift toward extrinsic motivators was thus justified; awareness alone could not bridge the gap between intention and action when systemic barriers remained high. Finally, the model identified social multipliers as the key to long-term sustainability. The significant influence of Social Norms ($w_{sn} \approx 0.35$) provided the mechanism for the ``Graduation Effect,'' allowing for the eventual reallocation of funds without a collapse in segregation rates \citep{Andre2021}. In conclusion, the DRL agent’s preference for localized enforcement and incentives was a sophisticated response to the socio-technical barriers in Philippine waste management, offering a data-driven roadmap toward a resilient compliance ecosystem \citep{Jimenez2025}.